\documentclass[11pt]{article}

\usepackage[margin=1in]{geometry}
\usepackage{amsmath,amssymb,mathtools}
\usepackage{booktabs}
\usepackage{enumitem}
\usepackage{hyperref}
\usepackage{listings}

\newcommand{\F}{\mathbb{F}_2}
\newcommand{\argminop}{\operatorname*{arg\,min}}
\newcommand{\rgb}{\{r,g,b\}}

\title{Near-Optimal Extensions of Concatenated MWPM for the 6.6.6 Color Code}
\author{Implementation specification}
\date{}

\begin{document}
\maketitle

\section{Scope}

This note defines two candidate-generation strategies for the concatenated
minimum-weight perfect matching (MWPM) decoder:
\begin{enumerate}
    \item \emph{cross-color relifting}, including pairwise relifting and an
    all-color, baseline-anchored relifting candidate;
    \item an \emph{X/Z-DEM prior perturbation ensemble}, in which the
    pre-color-decomposition detector error model is perturbed and then
    decomposed for each color.
\end{enumerate}

The two strategies should initially be evaluated independently. Candidate
generation may use a modified prior, but candidate selection must follow the
current common-prior scoring convention already implemented for
color-correlated decoding.

\section{Common original X/Z DEM}

Let the detector error model after the package's X/Z separation, but before
color decomposition, contain $m$ independent error mechanisms. Write
\[
    x\in\F^m
\]
for an error/correction hypothesis in this common original-DEM mechanism
ordering, and
\[
    H\in\F^{n\times m},\qquad Hx=s
\]
for its detector syndrome.

Let the original mechanism probabilities be
\[
    q=(q_1,\ldots,q_m),\qquad 0<q_i<1,
\]
with log-odds
\[
    \lambda_i=\log\frac{1-q_i}{q_i}.
    \tag{1}
\]

In the current repository this common model is represented by
\texttt{dem\_manager.dem\_xz} and \texttt{dem\_manager.probs\_xz}.

\section{Color decompositions}

For each target color $c\in\rgb$, define
\[
    H_{1,c}\in\F^{n_{1,c}\times m_{1,c}},\qquad
    H_{2,c}\in\F^{n_{2,c}\times m_{2,c}}.
\]
Let
\[
    A_{1,c}\in\F^{m_{1,c}\times m},\qquad
    A_{2,c}\in\F^{m_{2,c}\times m}
\]
be the provenance/error-map matrices from decomposed error mechanisms back to
original X/Z DEM mechanisms.

In \texttt{color-code-stim},
\begin{align*}
A_{1,c}
&\leftrightarrow
\texttt{dems\_decomposed[c].error\_map\_matrices[0]},\\
A_{2,c}
&\leftrightarrow
\texttt{dems\_decomposed[c].error\_map\_matrices[1]}.
\end{align*}

The Python implementation uses row-vector convention. Thus a native stage-2
prediction $v_c\in\F^{m_{2,c}}$ maps to original DEM ordering as
\[
    x_c=v_cA_{2,c}\pmod 2.
    \tag{2}
\]

\section{Baseline concatenated MWPM}

For each color $c$, ordinary concatenated MWPM produces a stage-1 hypothesis
\[
    u_c\in\F^{m_{1,c}}
\]
satisfying
\[
    H_{1,c}u_c=s_{1,c}.
    \tag{3}
\]

The decoder then forms the stage-2 syndrome
\[
    t_c(s,u_c)=
    \begin{pmatrix}
      s_c\\
      u_c
    \end{pmatrix},
    \tag{4}
\]
and runs stage-2 MWPM to obtain $v_c$, then maps it to $x_c$ using Eq.~(2).

Running this for $r,g,b$ requires six MWPM calls.

\section{Cross-color relifting}

\subsection{Required decomposition mode}

Cross-color relifting requires complete original-mechanism provenance and is
therefore unsupported when
\[
    \texttt{remove\_non\_edge\_like\_errors=True}.
\]
If cross-color relifting is enabled under this setting, the implementation must
raise an explicit error before decoding.

When \texttt{remove\_non\_edge\_like\_errors=False}, the implementation must
still verify that the resulting stage-1/stage-2 matrices are supported by the
current matching backend. If the full decomposition contains an unsupported
non-matchable column, fail explicitly rather than dropping it.

\subsection{Pairwise relifting}

Suppose baseline source color $d$ produces
\[
    x_d\in\F^m.
\]
For a different target color $c\neq d$, define
\[
    u_{c\leftarrow d}
    :=
    x_dA_{1,c}^{T}\pmod 2.
    \tag{5}
\]

Equation~(5) is a parity/XOR projection, not Boolean OR.

For a batch $X_d\in\F^{N\times m}$:
\begin{lstlisting}[language=Python,basicstyle=\ttfamily\small]
A1 = decomp_target.error_map_matrices[0]
u = ((X_d.astype(np.uint8) @ A1.T) % 2).astype(bool)
\end{lstlisting}

The required invariant is
\[
    H_{1,c}u_{c\leftarrow d}=s_{1,c}.
    \tag{6}
\]
This must be checked explicitly.

The pairwise relifted candidate skips target stage 1 and runs only target stage 2:
\[
    v_{c\leftarrow d}
    =
    D_{2,c}\left(t_c(s,u_{c\leftarrow d});q\right).
    \tag{7}
\]

\subsection{Why direct OR and unanchored two-source XOR are not used}

Let $d,e$ be the two colors different from target $c$.

Boolean OR,
\[
    u_{c\leftarrow d}\lor u_{c\leftarrow e},
\]
does not preserve the linear syndrome equation in general.

Direct XOR gives
\[
    H_{1,c}(u_{c\leftarrow d}\oplus u_{c\leftarrow e})
    =
    s_{1,c}\oplus s_{1,c}
    =
    0,
\]
so it is a zero-syndrome deformation, not a valid correction with syndrome
$s_{1,c}$.

\subsection{All-color baseline-anchored relifting}

Use baseline target-color solution $u_c$ as the anchor. Define
\[
    \delta_{c\leftarrow d}
    =
    u_c\oplus u_{c\leftarrow d},
    \qquad
    \delta_{c\leftarrow e}
    =
    u_c\oplus u_{c\leftarrow e}.
    \tag{8}
\]
Then
\[
    H_{1,c}\delta_{c\leftarrow d}
    =
    H_{1,c}\delta_{c\leftarrow e}
    =
    0.
    \tag{9}
\]

The all-color anchored hypothesis is
\begin{align}
    u_c^{\mathrm{all}}
    &=
    u_c
    \oplus
    \delta_{c\leftarrow d}
    \oplus
    \delta_{c\leftarrow e}
    \nonumber\\
    &=
    u_c
    \oplus
    u_{c\leftarrow d}
    \oplus
    u_{c\leftarrow e}.
    \tag{10}
\end{align}

Therefore
\[
    H_{1,c}u_c^{\mathrm{all}}=s_{1,c}.
    \tag{11}
\]

The corresponding physical candidate uses one additional target-color
stage-2 MWPM:
\[
    v_c^{\mathrm{all}}
    =
    D_{2,c}\left(t_c(s,u_c^{\mathrm{all}});q\right).
    \tag{12}
\]

This construction uses all three baseline branches: $u_c$ is the anchor and
the other two colors supply two independent zero-syndrome deviations.

\subsection{Candidate set and cost}

For one target color, the explored affine set is
\[
    u_c+
    \operatorname{span}_{\F}
    \{
      \delta_{c\leftarrow d},
      \delta_{c\leftarrow e}
    \}.
    \tag{13}
\]
The four explicit hypotheses are
\[
    u_c,\quad
    u_{c\leftarrow d},\quad
    u_{c\leftarrow e},\quad
    u_c^{\mathrm{all}}.
    \tag{14}
\]

Across three colors:
\[
    3\text{ baseline}
    +6\text{ pairwise relifts}
    +3\text{ all-color anchored relifts}
    =12
\]
physical candidates.

Without adaptive pruning, the full relifting family would require
\[
    6\text{ baseline}
    +6\text{ pairwise stage-2}
    +3\text{ all-color stage-2}
    =15
    \tag{15}
\]
MWPM calls per logical class.  This is the \emph{maximum} call count.  The
adaptive rules below must skip provably redundant instances, so the actual
per-shot count can be lower.

There is no relifting level parameter.


\section{Adaptive execution and early termination for relifting}

The relifting candidate family has a fixed \emph{logical} definition, but the
implementation should avoid executing MWPM instances that are guaranteed to
duplicate an already known result.

\subsection{Classify the three baseline physical corrections}

After ordinary concatenated MWPM, first map the three baseline final corrections
to the common original X/Z DEM mechanism ordering:
\[
    x_r,\;x_g,\;x_b\in\F^m.
\]

Define
\[
    N_{\rm uniq}
    :=
    \left|\{x_r,x_g,x_b\}\right|.
\]

For each shot, define the relifting run class
\[
    \rho_{\rm relift}
    =
    \begin{cases}
        0,&N_{\rm uniq}=1,\\
        1,&N_{\rm uniq}=2,\\
        2,&N_{\rm uniq}=3.
    \end{cases}
    \tag{23}
\]

Equality here means exact equality of the Boolean vectors in the original X/Z
DEM mechanism ordering.

This code is an execution-class label, not the number of MWPM calls.

\subsection{Case $\rho_{\rm relift}=0$: all three baseline corrections agree}

If
\[
    x_r=x_g=x_b,
\]
then no relifting instance is executed.  The ordinary baseline result is final,
because every source branch carries the same original-DEM hypothesis.

Thus the additional relifting cost is zero.

\subsection{Case $\rho_{\rm relift}=1$: exactly two baseline corrections agree}

Suppose, for example,
\[
    x_g=x_b\neq x_r.
\]

The two equal colors form one source-equivalence class.  Use the fixed color
order
\[
    r<g<b
\]
to choose a canonical representative whenever multiple equal source colors are
interchangeable.

For the example above:
\begin{itemize}
    \item target $r$: $g$ and $b$ are identical sources, so use only
    $r\leftarrow g$ because $g$ has the smaller color index;
    \item target $g$: the $b$ source is identical to the target baseline
    correction, so only $g\leftarrow r$ needs to be considered;
    \item target $b$: similarly, only $b\leftarrow r$ needs to be considered.
\end{itemize}

The all-color anchored candidates add no new stage-1 hypothesis in this case.
For example,
\[
    u_r^{\rm all}
    =
    u_r\oplus u_{r\leftarrow g}\oplus u_{r\leftarrow b}
    =
    u_r
\]
because the two relifted source hypotheses are equal.  For target $g$,
\[
    u_g^{\rm all}
    =
    u_g\oplus u_{g\leftarrow r}\oplus u_{g\leftarrow b}
    =
    u_{g\leftarrow r},
\]
because $x_b=x_g$ makes the $b$ relift equal to the target baseline hypothesis.
Therefore no separate all-color stage-2 call is required.

Before the stage-2 early-termination rule below, this case has at most three
additional stage-2 MWPM calls.

\subsection{Case $\rho_{\rm relift}=2$: all three baseline corrections differ}

If
\[
    x_r,\;x_g,\;x_b
\]
are pairwise distinct, generate the full relifting family:
six pairwise relifts and three all-color anchored relifts.

This gives at most nine additional stage-2 MWPM calls, before exact
stage-2-syndrome deduplication.

\subsection{Stage-2-syndrome early termination}

For a fixed target color $c$, define the baseline stage-2 syndrome
\[
    t_c^{(0)}:=t_c(s,u_c).
\]

For any proposed relifted stage-1 hypothesis $\tilde u$, construct the exact
stage-2 syndrome
\[
    \tilde t_c:=t_c(s,\tilde u)
\]
\emph{before} invoking stage-2 MWPM.

If
\[
    \tilde t_c=t_c^{(0)},
    \tag{24}
\]
then the target graph, target prior, and syndrome are exactly the same as the
baseline target-color stage-2 decoding problem.  Deterministic MWPM must
therefore return the same result.  This relifting instance terminates
immediately and reuses the baseline target-color candidate; stage-2 MWPM is not
called.

More generally, the implementation may safely cache already solved target-color
stage-2 syndromes:
\[
    (c,\tilde t_c)\longmapsto
    \text{stage-2 result}.
    \tag{25}
\]
If a later relifting candidate for the same target color produces a previously
solved syndrome, reuse that result rather than rerunning MWPM.  Equation~(24),
comparison against the baseline syndrome, is the required minimum behavior.

The equality test must compare the complete Boolean stage-2 detector vector
actually passed to PyMatching, including the target-color physical detector
part and the appended virtual detector/stage-1 part.

\subsection{Fixed logical candidate slots versus executed instances}

For analysis compatibility, the decoder may retain the canonical 12 candidate
slots
\[
    3\text{ baseline}
    +6\text{ single-source}
    +3\text{ all-color}
\]
even when some slots are aliases of an already computed candidate.

An aliased candidate should reuse the already computed:
\begin{itemize}
    \item mapped original correction,
    \item base-aligned native stage-2 correction,
    \item selection weight,
    \item observable prediction.
\end{itemize}

It must not count as an executed MWPM instance.

This preserves fixed candidate-array shapes while allowing dynamic runtime.

\subsection{Relifting execution sidecar}

The experiment layer should write a per-shot sidecar
\texttt{relift\_run.parquet}, following the exact shot-identity/schema
conventions used by the current \texttt{color\_correlated\_run.parquet} writer.

The required per-shot execution-class value is Eq.~(23):
\[
    0:\text{ all three baseline original-DEM corrections equal},
\]
\[
    1:\text{ exactly two baseline original-DEM corrections equal},
\]
\[
    2:\text{ all three baseline original-DEM corrections distinct}.
\]

The sidecar should additionally store the actual number of extra stage-2 MWPM
calls if the existing sidecar architecture admits such a field, because
Eq.~(24)--(25) can reduce the executed count below the class-level maximum.
The class code itself must remain defined only by baseline-correction equality.


\section{X/Z-DEM prior perturbation ensemble}

\subsection{Perturb the common pre-color-decomposition DEM}

Perturbation is applied to the common X/Z-separated DEM probability vector $q$,
not independently to the three color-decomposed stage-1 priors.

Let the total ensemble size be $M\ge1$. Member $m=0$ is exactly unperturbed:
\[
    q_i^{(0)}=q_i.
    \tag{16}
\]

For each member $m=1,\ldots,M-1$, sample
\[
    \xi_i^{(m)}\sim\operatorname{Uniform}[-1,1]
\]
independently and define
\[
    q_i^{(m)}
    =
    \operatorname{clip}
    \left(
      q_i(1+\alpha\xi_i^{(m)}),
      \epsilon,
      1-\epsilon
    \right),
    \qquad
    0\le\alpha\le1.
    \tag{17}
\]

Before clipping, this is equivalent to
\[
    q_i^{(m)}
    \sim
    \operatorname{Uniform}
    \left[
      (1-\alpha)q_i,
      (1+\alpha)q_i
    \right].
\]

For each member $m$, build one temporary X/Z DEM by replacing only the error
probabilities while preserving source ordering and targets. Then construct all
three color decompositions from that same temporary X/Z DEM:
\[
    \left(
      \mathcal D_{1,c}^{(m)},
      \mathcal D_{2,c}^{(m)}
    \right)
    =
    \operatorname{Decompose}_c
    \left(
      \mathcal D_{XZ}^{(m)}
    \right).
    \tag{18}
\]

Each color then runs a complete two-stage concatenated decoder:
\[
    u_c^{(m)}
    =
    D_{1,c}(s_{1,c};q^{(m)}),
\]
\[
    v_c^{(m)}
    =
    D_{2,c}(t_c(s,u_c^{(m)});q^{(m)}).
    \tag{19}
\]

Thus both induced stages are perturbed consistently from one common physical
prior.

The candidate count is $3M$ and the MWPM call count is $6M$ per logical class.

The perturbation table is fixed for the decoder instance/seed and reused across
shots, batches, and logical classes. It is not redrawn per shot.

\section{Candidate scoring under the unmodified base prior}

Candidate generation and candidate selection are distinct.

The current color-correlated implementation supports two common-prior scoring
bases. New strategies must use the same implementation path.

\subsection{Base stage-2 basis}

Let $p_{2,c}$ be the unmodified base stage-2 probabilities and
\[
    \lambda_{2,c,j}
    =
    \log\frac{1-p_{2,c,j}}{p_{2,c,j}}.
    \tag{20}
\]

A candidate must first be represented in the \emph{base} target-color stage-2
column ordering, denoted $\bar v_c$, and is scored as
\[
    C_{\mathrm{stage2}}(\bar v_c)
    =
    \langle\lambda_{2,c},\bar v_c\rangle.
    \tag{21}
\]

For a temporary decomposition, the native stage-2 order can differ from the
base order. The implementation must reuse the current
\texttt{align\_stage2\_to\_base} logic before applying Eq.~(21).

Relifting uses the base decomposition directly, so its native stage-2 vector is
already in base order.

\subsection{Original-DEM basis}

After mapping a candidate to original X/Z DEM ordering $x$, score
\[
    C_{\mathrm{org}}(x)
    =
    \langle\lambda,x\rangle,
    \tag{22}
\]
with $\lambda$ from Eq.~(1).

\subsection{Selection invariant}

All candidates in one decode call are compared under the same unmodified
base-prior basis. Temporary/generation weights from perturbed or conditioned
priors are diagnostic only.

For comparative decoding, minimize over candidate hypotheses within each
logical class first, then compare the logical-class minima. Logical gaps are
computed from those minima.

\section{Implementation invariants}

\begin{enumerate}
    \item With new options disabled, historical output is unchanged.
    \item Relifting with
    \texttt{remove\_non\_edge\_like\_errors=True} raises explicitly.
    \item Every pairwise and all-color relifted stage-1 hypothesis satisfies
    the target stage-1 syndrome equation.
    \item The all-color relift uses Eq.~(10), not OR and not unanchored
    two-source XOR.
    \item Perturbation acts on one common X/Z DEM and the same member prior
    generates all three color decompositions.
    \item Ensemble member $0$ is exactly the ordinary decoder.
    \item Base DEMs and base decompositions are never mutated.
    \item Temporary stage-2 outputs are aligned to base ordering before
    stage2-basis scoring.
    \item Selection weights use the current common-prior convention;
    generation weights never control final selection.
\end{enumerate}

\end{document}
